\section*{Capítulo \ref{ch:b2:thermal-radiation} --- Radiação térmica}

\begin{solution}{exo:b2:thermal-radiation:1}
$\lambda_{\max}$: \qty{0.50}{\micro m}, \qty{9.4}{\micro m}, \qty{1.04}{\micro m},
\qty{1.06}{mm}, \qty{38}{\micro m}, \qty{3.2}{\micro m}; $\sigma T^4$: \qty{6.4e7}{},
\qty{520}{}, \qty{3.5e6}{}, \qty{3.1e-6}{}, \qty{2.0}{}, \qty{3.7e4}{W/m^2}. Visíveis:
o Sol e o filamento; o atiçador brilha em vermelho-escuro pela cauda de onda
curta, um décimo de milésimo de sua potência.
\end{solution}

\begin{solution}{exo:b2:thermal-radiation:2}
$S = \sigma T_\odot^4(R_\odot/d)^2 = \qty{1360}{W/m^2}$; $L = 4\pi d^2S =
\qty{3.8e26}{W}$; $\pi R^2S = \qty{1.7e17}{W}$; \qty{2e7}{W} por pessoa --- nove
mil vezes o consumo mundial.
\end{solution}

\begin{solution}{exo:b2:thermal-radiation:3}
Emitidos \qty{830}{W}, absorvidos \qty{700}{W}, líquido \qty{130}{W}; $h_{\text{r}} =
4\varepsilon\sigma T_0^3 = \qty{5.6}{W/m^2/K}$. Mais que o metabolismo:
as roupas reduzem isso; paredes a \qty{15}{\celsius} acrescentam uns \qty{50}{W} de perda
--- sente-se frio em ar morno entre paredes frias.
\end{solution}

\begin{solution}{exo:b2:thermal-radiation:4}
$A = P/\varepsilon\sigma T^4 = \qty{49}{mm^2}$; $\ell = A/\pi d = \qty{31}{cm}$;
\qty{4.8}{W} de luz; um LED de \qty{16}{W}.
\end{solution}

\begin{solution}{exo:b2:thermal-radiation:5}
(a) $\eu^x - 1 \approx x$: $u_\nu = (8\pi\nu^2/c^3)\,k_BT$ --- $k_BT$ por modo.
(b) $\int\nu^2\dd\nu$ diverge. (c) $\nu = k_BT/h = \qty{6e12}{Hz}$ (\qty{50}{\micro m}):
o rádio e as micro-ondas a \qty{300}{K} estão bem dentro do regime clássico (o
ruído de um resistor vale $k_BT$ por unidade de banda). (d) A probabilidade de
encontrar o quantum $h\nu$ excitado a $T$: um fator de Boltzmann.
\end{solution}

\begin{solution}{exo:b2:thermal-radiation:6}
(a) $x = 4.965$. (b) $hc/4.965k_B = \qty{2.90e-3}{m.K}$. (c) $x^3/(\eu^x - 1)$
tem máximo onde $3(1 - \eu^{-x}) = x$: $2.82$; as densidades por unidade de frequência
e por unidade de comprimento de onda são funções diferentes ($\dd\nu = c\,\dd\lambda/\lambda^2$)
e têm máximo em lugares diferentes. (d) \qty{10000}{K}; branco azulado.
\end{solution}

\begin{solution}{exo:b2:thermal-radiation:7}
(a) $u = (8\pi^5k_B^4/15h^3c^3)T^4$, $\sigma = cu/4T^4 = \qty{5.67e-8}{W/m^2/K^4}$.
(b) A isotropia e a projeção em $\cos\theta$: $\int\cos\theta\,\dd\Omega/4\pi = 1/4$
sobre um hemisfério. (c) $n = 8\pi \times 2.404\,(k_BT/hc)^3 = 60.4\,(k_BT/hc)^3$;
energia média $(\pi^4/15)/2.404 = 2.70\,k_BT$. (d) $5 \times 10^8$ e $400$
por centímetro cúbico.
\end{solution}

\begin{solution}{exo:b2:thermal-radiation:8}
(a) $\sigma(T_1^4 - T_2^4)$. (b) $T_{\text{s}}^4 = (T_1^4 + T_2^4)/2$; a metade. (c)
$1/(n + 1)$. (d) Somando as reflexões, $\sigma(T_1^4 - T_2^4)/(2/\varepsilon -
1)$; $\varepsilon = 0.05$: $\qty{460}{W/m^2}/39 = \qty{12}{W/m^2}$ --- uma garrafa de
\qty{0.05}{m^2} perde \qty{0.6}{W}, evaporando um quarto de quilograma de
nitrogênio por dia.
\end{solution}

\begin{solution}{exo:b2:thermal-radiation:9}
(a) $\varepsilon\sigma(T^4 - 250^4) = h(278 - T)$: $T \approx \qty{266}{K}$. (b)
\qty{277}{K}. (c) Um $h$ maior prende a superfície ao ar. (d) A copa
a \qty{278}{K} substitui o céu a \qty{250}{K}.
\end{solution}

\begin{solution}{exo:b2:thermal-radiation:10}
(a) \qty{231}{K}, \qty{255}{K}, \qty{210}{K}. (b) \qty{506}{K}, \qty{33}{K}, \qty{5}{K}.
(c) Cada camada opaca irradia para cima e para baixo; os balanços dão
$T_k^4 = (k + 1)T_{\text{e}}^4$ a partir do topo: $n \approx 100$ para Vênus. (d)
$((1 - A)S/\sigma)^{1/4} = \qty{381}{K}$; sem atmosfera nada leva
calor ao lado noturno, e o regolito armazena pouco: a superfície
irradia o calor do dia em cerca de um dia e se estabiliza perto de
\qty{100}{K}.
\end{solution}

\begin{solution}{exo:b2:thermal-radiation:11}
(a) $4\sigma T_{\text{e}}^3 = \qty{3.8}{W/m^2/K}$. (b) \qty{1.0}{K}. (c) $C =
\qty{2.1e8}{J/m^2/K}$, $\tau = C/4\sigma T_{\text{e}}^3 \approx \qty{5.6e7}{s}$, dois
anos. (d) Ar mais quente retém mais vapor de água (um gás de \omterm{prop:b2:thermal-radiation:greenhouse}{efeito estufa}) e
derrete gelo (albedo menor): ambos amplificam o aquecimento inicial.
\end{solution}

\begin{solution}{exo:b2:thermal-radiation:12}
(a)
\[
\frac{M_1}{M_2} = \Big(\frac{\lambda_2}{\lambda_1}\Big)^5\exp\Big[-\frac{hc}{k_BT}\Big(\frac{1}{\lambda_1} - \frac{1}{\lambda_2}\Big)\Big] :
\]
$\varepsilon$ se cancela. (b) $T = \qty{14388}{\micro m.K} \times 0.427/(5\ln 1.385 +
\ln 10) = \qty{1570}{K}$. (c) Recebido $\propto 0.1 \times 350 + 0.9 \times 293 =
299$; indicado $(299 - 0.05 \times 293)/0.95 \approx \qty{299}{K}$: ela erra os
\qty{350}{K}. (d) Sua \omterm{prop:b2:thermal-radiation:kirchhoff}{emissividade} é justamente a suposta.
\end{solution}

\begin{solution}{pb:b2:thermal-radiation:1}
\textbf{1.} \qty{5800}{K}.

\textbf{2.} \qty{6.4e7}{W/m^2}; \qty{3.9e26}{W}.

\textbf{3.} $L/4\pi d^2 = \qty{1.37e3}{W/m^2}$; $u = S/c = \qty{4.6e-6}{J/m^3}$;
$P = S/c = \qty{4.6e-6}{Pa}$ se absorvida, o dobro se refletida.

\textbf{4.} $2.7k_BT = \qty{2.2e-19}{J}$ (\qty{1.35}{eV}); $6 \times 10^{21}$
fótons por segundo e por metro quadrado.

\textbf{5.} $3.9 \times 10^{26} \times 1.4 \times 10^{17} = \qty{5.5e43}{J}$; $Mc^2 =
\qty{1.8e47}{J}$: três décimos de milésimo de sua massa.

\textbf{6.} \qty{1.2e17}{W}; $T_{\text{e}} = \qty{255}{K}$.

\textbf{7.} Interceptada em $\pi R^2$, emitida de $4\pi R^2$ graças à
rotação e aos ventos; só o ponto subsolar: $((1 - A)S/\sigma)^{1/4} =
\qty{360}{K}$.

\textbf{8.} \qty{381}{K}; à noite o calor armazenado ($\sim C\Delta T$) escoa por
radiação com a constante de tempo $C/4\sigma T^3 \approx \qty{1}{dia}$, curta
frente à quinzena de escuridão: o solo cai até onde sua
radiação iguala o filete que vem de baixo, \qty{100}{K}.

\textbf{9.} \qty{10}{\micro m}, vinte vezes o do Sol.

\textbf{10.} $(0.65/0.70)^{1/4} = 0.982$: \qty{4.7}{K} menos, \qty{250}{K}.

\textbf{11.} Espaço: $\sigma T_{\text{a}}^4 = \sigma T_{\text{e}}^4$; camada: $\sigma
T_{\text{s}}^4 = 2\sigma T_{\text{a}}^4$; superfície: $\sigma T_{\text{s}}^4 = \sigma
T_{\text{e}}^4 + \sigma T_{\text{a}}^4$. Logo $T_{\text{s}} = 2^{1/4}T_{\text{e}} = \qty{303}{K}$.

\textbf{12.} $T_{\text{s}}^4 = T_{\text{e}}^4/(1 - \varepsilon/2)$; $(288/255)^4 = 1.63$
dá $\varepsilon = 0.77$.

\textbf{13.} O vapor de água, o dióxido de carbono, o metano e o ozônio absorvem no
infravermelho (bandas de vibração: H$_2$O em torno de \qty{6}{\micro m} e além de
\qty{17}{\micro m}, CO$_2$ em \qty{15}{\micro m}); N$_2$ e O$_2$ não (sem
dipolo); a janela vai de \qty{8}{} a \qty{13}{\micro m}.

\textbf{14.} $T_{\text{a}}^4 = T_{\text{s}}^4/2$: $\varepsilon\sigma T_{\text{a}}^4 \approx
\qty{150}{W/m^2}$, contra \qty{238}{W/m^2} de luz solar absorvida (uma
atmosfera de várias camadas devolve mais, uns \qty{330}{W/m^2}).

\textbf{15.} $3.7/3.8 \approx \qty{1}{K}$.

\textbf{16.} \qty{2.1e8}{J/m^2/K}; $\tau \approx \qty{5.6e7}{s}$, dois anos.

\textbf{17.} Elas refletem a luz do Sol e irradiam infravermelho para baixo; à noite
só o segundo age: as noites nubladas são mais quentes.

\textbf{18.} \qty{49}{mm^2}; \qty{31}{cm}; enrolado para caber e para reduzir
as perdas convectivas para o gás de enchimento.

\textbf{19.} \qty{1.04}{\micro m}; \qty{4.8}{W}; $8\%$.

\textbf{20.} $R = V^2/P = \qty{880}{\ohm}$; frio \qty{88}{\ohm}: \qty{2.6}{A}
em vez de \qty{0.26}{A} --- as lâmpadas queimam ao serem ligadas.

\textbf{21.} $49 \times (2800/3200)^4 = \qty{29}{mm^2}$; \qty{8.4}{W} visíveis; oito
duplicações da evaporação, uma vida $250$ vezes mais curta --- horas.

\textbf{22.} O tungstênio evaporado se liga ao halogênio e é devolvido
ao filamento quente em vez de enegrecer o vidro: um filamento mais quente
com vida normal, num pequeno bulbo de quartzo.

\textbf{23.} \qty{55}{W} de infravermelho e de condução pelo gás; o
vidro absorve parte do infravermelho.

\textbf{24.} $C = \qty{2.8e-3}{J/K}$; com $P \propto T^4$, o tempo para reduzir $T$ à metade
vale $7CT_0/3P \approx \qty{0.3}{s}$ --- e o brilho segue a rede
com um atraso de \qty{0.04}{s}, razão pela qual ele quase não cintila.

\textbf{25.} A temperatura fixa a forma e o pico do espectro
($\lambda_{\max}T$), e o total ($\sigma T^4$); filamento, Sol e Terra
são três \omterm{def:b2:thermal-radiation:blackbody}{corpos negros}, a \qty{2800}{}, \qty{5800}{} e \qty{288}{K}, cada um
equilibrando o que recebe contra $\sigma T^4$.
\end{solution}
